\begin{tikzpicture}[
    font=\small,
    flow/.style={
        draw=black!80,
        -Latex,
        line width=0.8pt,
        shorten <=1.5pt,
        shorten >=1.5pt
    },
    inputbox/.style={
        draw=black!75,
        rounded corners=4pt,
        fill=orange!8,
        text width=4.0cm,
        minimum height=4.9cm,
        align=center,
        line width=0.8pt,
        inner xsep=7pt,
        inner ysep=7pt
    },
    branchbox/.style={
        draw=black!75,
        rounded corners=4pt,
        text width=5.2cm,
        minimum height=2.85cm,
        align=center,
        line width=0.8pt,
        inner xsep=7pt,
        inner ysep=6pt
    },
    positivebox/.style={
        branchbox,
        fill=green!9
    },
    negativebox/.style={
        branchbox,
        fill=red!6
    },
    resultbox/.style={
        draw=black!75,
        rounded corners=4pt,
        text width=4.2cm,
        minimum height=2.35cm,
        align=center,
        line width=0.8pt,
        inner xsep=7pt,
        inner ysep=6pt
    },
    positiveresult/.style={
        resultbox,
        fill=green!6
    },
    negativeresult/.style={
        resultbox,
        fill=red!4
    },
    principlebox/.style={
        draw=black!65,
        rounded corners=4pt,
        fill=gray!6,
        text width=13.3cm,
        minimum height=1.2cm,
        align=center,
        line width=0.7pt,
        inner xsep=8pt,
        inner ysep=6pt
    },
    arrowlabel/.style={
        font=\footnotesize,
        text=black!80,
        fill=white,
        inner xsep=2pt,
        inner ysep=1pt
    },
    branchlabel/.style={
        font=\footnotesize,
        text=black!80,
        fill=white,
        inner xsep=2pt,
        inner ysep=1pt
    }
]

% 设置安全边界，避免导出时裁切
\path[use as bounding box]
    (-0.2,-0.35) rectangle (18.2,8.2);

% ==================================================
% 左侧：当前状态、策略与动作采样
% ==================================================

\node[inputbox] (input) at (2.2,4.75) {
    {\bfseries 当前策略与动作}\\[6pt]

    {\bfseries 当前棋盘状态 $s_t$}\\[4pt]
    候选位置：$B=(8,8)$\\[4pt]
    当前策略概率：\\[2pt]
    $\displaystyle
    \pi_{\theta}(B\mid s_t)=0.40
    $\\[9pt]

    {\footnotesize 按当前策略采样}\\[-1pt]
    $\Downarrow$\\[3pt]

    $\displaystyle
    a_t=B=(8,8)
    $\\[7pt]

    {\footnotesize
    随后观察该动作产生的实际回报}
};

% ==================================================
% 中间：正优势分支
% ==================================================

\node[positivebox] (positive) at (9.0,6.40) {
    {\bfseries 正优势（$+$）：回报高于基线}\\[6pt]

    实际回报：$G_t=8$\\[3pt]
    状态基线：$b(s_t)=6$\\[5pt]

    优势估计：$\widehat{A}_t=G_t-b(s_t)$\\[2pt]
    $\displaystyle
    =8-6=2>0
    $\\[5pt]

    {\footnotesize
    该动作表现优于当前局面的平均水平}
};

% ==================================================
% 中间：负优势分支
% ==================================================

\node[negativebox] (negative) at (9.0,2.10) {
    {\bfseries 负优势（$-$）：回报低于基线}\\[6pt]

    实际回报：$G_t=3$\\[3pt]
    状态基线：$b(s_t)=6$\\[5pt]

    优势估计：$\widehat{A}_t=G_t-b(s_t)$\\[2pt]
    $\displaystyle
    =3-6=-3<0
    $\\[5pt]

    {\footnotesize
    该动作表现低于当前局面的平均水平}
};

% ==================================================
% 右侧：正优势对应的策略更新结果
% ==================================================

\node[positiveresult] (positive-result) at (15.6,6.40) {
    {\bfseries 动作概率提高（$\uparrow$）}\\[7pt]

    更新前：\\[-1pt]
    $\displaystyle
    \pi_{\theta}(B\mid s_t)=0.40
    $\\[5pt]

    更新后：\\[-1pt]
    $\displaystyle
    \pi_{\theta}(B\mid s_t)=0.48
    $\\[6pt]

    {\footnotesize
    后续更容易选择动作 $B$}
};

% ==================================================
% 右侧：负优势对应的策略更新结果
% ==================================================

\node[negativeresult] (negative-result) at (15.6,2.10) {
    {\bfseries 动作概率降低（$\downarrow$）}\\[7pt]

    更新前：\\[-1pt]
    $\displaystyle
    \pi_{\theta}(B\mid s_t)=0.40
    $\\[5pt]

    更新后：\\[-1pt]
    $\displaystyle
    \pi_{\theta}(B\mid s_t)=0.32
    $\\[6pt]

    {\footnotesize
    后续选择动作 $B$ 的概率下降}
};

% ==================================================
% 左侧动作节点到正负分支的公共分叉
% ==================================================

\coordinate (split) at (5.35,4.75);
\coordinate (positive-join) at (5.35,6.40);
\coordinate (negative-join) at (5.35,3.10);

% 从左侧动作模块引出主干
\draw[flow]
    (input.east)
    --
    (split);

% 公共分叉点
\fill[black!80]
    (split) circle (1.4pt);

% 正优势分支
\draw[flow]
    (split)
    --
    (positive-join)
    --
    node[midway, above=3pt, branchlabel]
    {$G_t>b(s_t)$}
    (positive.west);

% 负优势分支
\draw[flow]
    (split)
    --
    (negative-join)
    --
    node[midway, below=3pt, branchlabel]
    {$G_t<b(s_t)$}
    (negative.west);

% ==================================================
% 优势分支到概率更新结果
% ==================================================

\draw[flow]
    (positive.east)
    --
    node[midway, above=4pt, arrowlabel]
    {提高选择概率}
    (positive-result.west);

\draw[flow]
    (negative.east)
    --
    node[midway, above=4pt, arrowlabel]
    {降低选择概率}
    (negative-result.west);

% ==================================================
% 底部：统一的基线与策略更新原则
% ==================================================

% \node[principlebox] (principle) at (9.0,-0.7) {
%     {\bfseries 基线与策略更新原则}\\[5pt]
%     常用状态基线为
%     $\displaystyle b(s_t)=V^{\pi}(s_t)$。\\[3pt]
%     策略更新关注的不是回报的绝对大小，
%     而是动作结果相对于当前平均水平的好坏。
% };

\end{tikzpicture}